\section{Related work, positioned}

Table~\ref{tab:optimisations} lists what we adopted and what we rejected, with the evidence that
rejected it. The related systems differ from ours in mechanism and in claim, and it is worth being
precise about both.

\paragraph{Declared-action caches} (Bazel, Vesta, Nix) assume the action states its inputs. We do
not compete with them; we measure what it costs to operate without that assumption. A system with
declarations should use them.

\paragraph{Compiler caches} (\texttt{ccache}, \texttt{sccache}) are specialists and win their home
ground by a wide margin, as Section~\ref{sec:negative} reports. Our contribution is not a
compiler-cache replacement and must not be marketed as one.

\paragraph{\firebuild{}} is the closest system: an automatic cache for unmodified commands, using
library interposition rather than tracing, with its own validation of recorded inputs. Our
comparison is against two of its versions on the same host. The differences we can defend are
narrow and specific: applicability for timestamp-setting producers, behaviour on observations that
cannot be reproduced (it replays clock-dependent output under its default configuration, and its
strict configuration declines to cache it), and the handling of a corrupted cache object
(it serves the corrupted bytes, we refuse and re-execute). Each is a configuration-dependent
statement, not a ranking.

\paragraph{Record/replay systems} such as \texttt{rr} reproduce an execution exactly by logging its
non-deterministic inputs. They provide a strictly stronger guarantee and pay a correspondingly
higher recording and storage cost. We do not claim their properties, and a reader who needs an
exact replay should use one.

\paragraph{Deterministic build literature} motivates both sides of our trade-off: hermeticity makes
caching easy, and non-hermetic commands are exactly where automatic caches must be conservative or
wrong~\cite{reproduciblebuilds}.

\section{Conclusion}

We set out to answer what a deliberately conservative automatic command cache costs and what it
buys, and the answer is more precise than a speedup.

\emph{What it buys.} Across a frozen 59-workload census and every measurement in this paper, no
reuse produced an observation that differed from a fresh execution, and the mechanism that makes
that true is auditable: a guard per recorded observation, a single validation pass adjacent to the
restore, a typestate gate that makes unverified bytes unexposable, and a re-hash of every cached
object before exposure. Within the workload class the system targets --- repeated generation of
stable-input artifacts under a controlled contract --- the saving is large and repeatable:
$5.22\times$ and $80.8\%$ less cumulative time on a 25-invocation \texttt{bzip2} session, on a
workload the closest competing system declines to cache at all.

\emph{What it costs.} Eligibility is $4.1\%$ natively and $24.5$--$38.8\%$ under an audited
contract, so most commands simply never hit. Warm hits cost milliseconds, so cheap commands lose.
Verified replay costs tens to hundreds of milliseconds at large artifacts, so a cache with a
cheaper validation story beats us there. Compilation belongs to compiler caches. The enforced tier
completes its lifecycle for read-only commands but cannot promote the producers that motivate the
system. And a mutation after validation still yields a stale hit, which no economical userspace
mechanism prevents.

\emph{What we learned about conservative execution caching.} Refusal is not free and not
self-justifying: it must be measured in eligibility, in latency and in applicability, or a system
cannot tell whether it is being careful or merely slow. The useful engineering is not the tracing
--- that is composition --- but the accounting: what is observed, what is gateable, what is
disclosed, and what is retained as a counterexample instead of being tuned away. A cache that can
say exactly which of its results are trustworthy, and why, is more valuable than one that is faster
and cannot.

\section{Artifact availability}

The artefact is the repository, and every number in this paper is regenerated from committed
evidence rather than transcribed.

\begin{itemize}
  \item \textbf{Repository}: \url{https://github.com/Satwik-P28/tracejit}.
  \item \textbf{Source revisions}: the research evidence in this paper was produced on
  \texttt{b81242d65e8a02bfe4459506a44b2eca566c5d89} (PR \#45) with \texttt{main} at
  \texttt{68ba4067afd4391c90af9a64ce0b43397438d504}; the Phase AF enforcement and artifact-scale
  evidence in this paper was produced on branch \texttt{research/phase-af} from that head. The
  proposed release is \texttt{0.3.0-alpha.1} and is not published.
  \item \textbf{Tables}: generated by
  \texttt{python3 benchmarks/phase\_af/build\_tables.py --emit paper/tables}; the generator reads
  \texttt{benchmarks/results/**} and fails when a document's headline disagrees with the rows
  (\texttt{--check}). \texttt{benchmarks/phase\_af/test\_tables.py} fails when a table loses a row,
  gains an unknown cell, or the census stops reconciling.
  \item \textbf{Raw evidence}: \texttt{benchmarks/results/phaseaa1/} (census),
  \texttt{phaseae/} (competitor, integrity, race, replay, environments),
  \texttt{phaseac/} (metadata and stale-archive counterexamples),
  \texttt{phaseaf/} (enforcement lifecycle, artifact scale).
  \item \textbf{Reproduction}: \texttt{docs/REPRODUCTION\_PACKAGE.md} specifies a small path (build,
  guarded reuse, mutation, corrupt-object recovery, a compression workload, a short correctness
  suite) and a full path (both \firebuild{} versions, the competitive matrix, the census, the
  artifact study, adversarial suites, environment qualification, table regeneration), with pinned
  versions, expected runtime, storage and capability-dependent skips.
  \item \textbf{Host fingerprint} for the primary timing host:
  \texttt{benchmarks/results/phaseaa1/environment.json}.
\end{itemize}

\begin{table}[t]
\centering
\small
\caption{Candidates accepted into the product and candidates rejected with the measurement that rejected them.}
\label{tab:optimisations}
\begin{tabular}{lrr}
\toprule
\textbf{Candidate} & \textbf{Status} & \textbf{Evidence} \\
\midrule
Drop the redundant first guard pass before loading cached streams & accepted & docs/PHASE\_X\_REPORT.md; one pass remains, immediately before restore; guard-pass regression suite in tracejit-core \\
Skip the durable restore intent for records with no workspace outputs & accepted & docs/PHASE\_AB\_REPORT.md; paired A/B on cut, gzip -6, sed with CIs excluding zero \\
Drop the pre-delete/pre-rename input snapshot when the same execution recreates the path & accepted & docs/PHASE\_W\_REPORT.md; regression superseded\_delete\_does\_not\_pin\_the\_next\_run\_to\_the\_old\_bytes \\
Report the record's determinism class on the reuse path & accepted & docs/PHASE\_AF\_REPORT.md; linux\_integration::enforced\_execution\_is\_proven\_reused\_and\_invalidated \\
Promote F\_GETFD on inherited descriptors to a modelled no-op & rejected & docs/PHASE\_AD\_REPORT.md §14: sound but inert, 0 of 48 census cells changed; unit test pins the refusal \\
Model F\_GETFL/F\_SETFL on inherited or internal pipes & rejected & docs/PHASE\_AD\_REPORT.md §14: shared, externally mutable status flags \\
Model tar's realtime clock read & rejected & docs/PHASE\_AD\_REPORT.md §14: no clock contract exists; --mtime does not remove the read \\
Relax metadata guards with a metadata-preserving clone & rejected & docs/PHASE\_AC\_REPORT.md: unjustifiable field precision, reproduced divergence \\
Snapshot inputs to close the validation-to-restore race & rejected & docs/TOCTOU\_CORRECTNESS\_BOUNDARY.md: redefines observable state and doubles hit cost \\
Cooperative flock mode for reuse & rejected & docs/TOCTOU\_CORRECTNESS\_BOUNDARY.md: not a guarantee against an uncooperative writer \\
Change default admission from auto to always & rejected & docs/PHASE\_AD\_REPORT.md §7: 3.4-4.2\% on one session, unproven, costs one-shot tracing \\
Fused hash-and-write, streaming cached bytes, mmap replay & rejected & docs/PHASE\_AE\_REPORT.md §13: each exposes unverified bytes or adds a failure mode \\
Cache guard results between runs (metadata-only revalidation) & rejected & docs/PHASE\_AE\_REPORT.md §13: reintroduces the same-timestamp defect \\
Parallel guard validation across 8 threads & rejected & docs/PHASE\_W\_REPORT.md: sequential guard walk dominates; 111.7 to 110.6 ms median, reverted \\
Grant the parent directory of a declared output REMOVE\_FILE/MAKE\_REG so gzip can be enforced & rejected & docs/PHASE\_AF\_REPORT.md §5: does not help (the refusal is an unmodelled EEXIST write) and widens the policy \\
Grant parent-directory replacement rights to make producers enforceable & rejected & docs/PHASE\_AF\_REPORT.md §5: the blocking effect is the command's own failed O\_EXCL, not the directory grant \\
\bottomrule
\end{tabular}
\end{table}

